\documentclass[10pt, a4paper]{article}

% --- Packages ---
\usepackage[utf8]{inputenc}
\usepackage[T1]{fontenc}
\usepackage[margin=0.60in, top=0.72in, bottom=0.72in]{geometry}
\usepackage{xcolor}
\usepackage{titlesec}
\usepackage{booktabs}
\usepackage{tabularx}
\usepackage{longtable}
\usepackage{enumitem}
\usepackage{amsmath, amssymb}
\usepackage{fancyhdr}
\usepackage{microtype}
\usepackage{multicol}
\usepackage{mdframed}
\usepackage[hidelinks]{hyperref}

% --- Monochromatic Color Palette ---
\definecolor{monodark}{HTML}{111111}
\definecolor{monomedium}{HTML}{3A3A3A}
\definecolor{monolight}{HTML}{666666}
\definecolor{monoline}{HTML}{CCCCCC}
\definecolor{monoshade}{HTML}{F0F0F0}
\definecolor{monobox}{HTML}{E8E8E8}

% --- Typography & Spacing ---
\setlist[itemize]{noitemsep, topsep=2pt, leftmargin=14pt}
\setlist[enumerate]{noitemsep, topsep=2pt, leftmargin=14pt}
\emergencystretch=4em

% --- Code Box Style ---
\newmdenv[
  backgroundcolor=monoshade,
  linecolor=monoline,
  linewidth=0.5pt,
  topline=true,
  bottomline=true,
  leftline=true,
  rightline=true,
  innerleftmargin=6pt,
  innerrightmargin=6pt,
  innertopmargin=4pt,
  innerbottommargin=4pt,
  skipabove=4pt,
  skipbelow=4pt,
]{codebox}

% --- Section Headers ---
\titleformat{\section}
  {\color{monodark}\large\bfseries}
  {\thesection.}{0.5em}{}
  [\color{monoline}\titlerule]

\titleformat{\subsection}
  {\color{monomedium}\normalsize\bfseries}
  {\thesubsection}{0.5em}{}

\titleformat{\subsubsection}
  {\color{monolight}\small\bfseries\itshape}
  {}{0em}{}

\titlespacing*{\section}{0pt}{10pt}{4pt}
\titlespacing*{\subsection}{0pt}{7pt}{2pt}
\titlespacing*{\subsubsection}{0pt}{5pt}{1pt}

% --- Header and Footer ---
\pagestyle{fancy}
\fancyhf{}
\renewcommand{\headrulewidth}{0.4pt}
\renewcommand{\footrulewidth}{0.4pt}
\rfoot{\small\color{monolight} Page \thepage}

% =============================================================
\begin{document}
% =============================================================

% ---- TITLE BLOCK ----
\begin{center}
  {\Huge\bfseries\color{monodark} Smart Flow}\\[5pt]
  {\large\itshape\color{monomedium}
    AI Powered Traffic Signal Control via Computer Vision \& Deep Reinforcement Learning
  }\\[9pt]
  \noindent\rule{\linewidth}{0.4pt}\\[6pt]
\end{center}

% ---- ABSTRACT ----
\noindent\textbf{Abstract.}\quad
Smart Flow is an AI driven adaptive traffic signal control system developed for
Smart India Hackathon 2024. It combines a computer vision pipeline (lane
detection via Hough Transform, background subtraction, contour-based vehicle
counting, and YOLOv8 object detection) with a Deep Q-Network (DQN) trained
inside the SUMO traffic microsimulator. The system observes live vehicle
density per lane, then dynamically adjusts traffic-light phase durations to
minimise cumulative waiting time across all controlled intersections. This
document provides a full technical specification: modular package architecture, CV
algorithms, RL agent design, SUMO integration, map data, trained model
artefacts, data-flow diagrams, engineering decisions, refactored project structure,
and a future roadmap.

\vspace{6pt}
\noindent\rule{\linewidth}{0.4pt}
\vspace{4pt}

% ---- SWITCH TO TWO-COLUMN BODY ----
\raggedcolumns
\begin{multicols}{2}

% ==========================================
\section{Introduction \& Motivation}
% ==========================================

Urban intersections in India's growing cities have become critical
bottlenecks. Average commuters lose 2--3 hours daily; conventional
timer-based signals cannot adapt to dynamic patterns; idling vehicles emit
significant CO\textsubscript{2}. Manual traffic management is inconsistent
and resource-intensive.

Smart Flow addresses every one of these gaps by (1) \textit{seeing} real
traffic conditions through computer vision, (2) \textit{learning} an optimal
switching policy through reinforcement learning, and (3) \textit{acting} by
dynamically adjusting signal phase durations in real time.

\subsection{Project Goals}
\begin{enumerate}
  \item Automatically detect the number of vehicles in each lane from a
        CCTV-quality video feed, with no specialised hardware.
  \item Train a DQN agent inside the SUMO microsimulator to discover a
        policy that minimises lane waiting time.
  \item Generalise across multiple real-world intersection topologies
        (single intersection, city block, King Circle roundabout).
  \item Export the trained model so it can be deployed on edge hardware
        via ONNX $\to$ OpenVINO inference.
\end{enumerate}

% ==========================================
\section{Complete Technology Stack}
% ==========================================

\begin{center}
\small\color{monodark}
\begin{tabular}{@{}p{1.7cm} p{4.5cm}@{}}
\toprule
\textbf{Layer} & \textbf{Technology} \\
\midrule
Language     & Python 3.10+ \\
CV Library   & OpenCV (\texttt{cv2}) 4.x \\
Detection    & YOLOv8n (\texttt{ultralytics}) \\
RL Framework & Stable-Baselines3 (DQN) \\
RL Env       & Gymnasium (\texttt{gymnasium}) \\
Simulator    & SUMO + TraCI API \\
Map Source   & OpenStreetMap $\to$ \texttt{netconvert} \\
Numerics     & NumPy, Matplotlib \\
Notebooks    & Jupyter (exploration only) \\
Model Format & SB3 \texttt{.zip} $\to$ ONNX $\to$ OpenVINO IR \\
\bottomrule
\end{tabular}
\end{center}

% ==========================================
\section{Repository File Structure}
% ==========================================

The repository has been restructured into clean, modular Python packages
and dedicated asset directories.

\subsection{Layout Mapping}

\begin{center}
\scriptsize
\begin{tabular}{@{}p{3.5cm} p{4.3cm}@{}}
\toprule
\textbf{File path} & \textbf{Description} \\
\midrule
\texttt{cv/utils.py}
  & contour detection \& cropping functions \\
\texttt{cv/lane\_detector.py}
  & Core lane landmark API \\
\texttt{cv/background\_contour.py}
  & Noise baseline estimation \\
\texttt{cv/vehicle\_counter.py}
  & Real time contour vehicle counter \\
\texttt{cv/yolo\_counter.py}
  & YOLOv8 object detection counter \\
\texttt{cv/lane\_debug\_image.py}
  & Still image lane visualizer \\
\texttt{cv/lane\_debug\_video.py}
  & Video lane visualizer \\
\texttt{cv/hough\_explore\_image.py}
  & Hough slope exploration (image) \\
\texttt{cv/hough\_explore\_video.py}
  & Hough slope exploration (video) \\
\texttt{rl/train\_dqn.py}
  & DQN training entry point \\
\texttt{rl/test\_dqn.py}
  & SUMO-GUI evaluation script \\
\texttt{rl/saved\_models/}
  & Trained model \texttt{.zip} archives \\
\texttt{sumo/maps/}
  & \texttt{.net.xml} \& \texttt{.rou.xml} files \\
\texttt{notebooks/}
  & Jupyter notebooks \\
\bottomrule
\end{tabular}
\end{center}

% ==========================================
\section{High Level System Architecture}
% ==========================================

The system consists of two decoupled subsystems sharing the lane density
signal:

\subsubsection{Subsystem A: Computer Vision Pipeline (\texttt{cv/})}
Processes CCTV video feeds to generate real-time vehicle counts per lane.
Built with OpenCV and NumPy, executing standalone without simulator overhead.

\subsubsection{Subsystem B: RL Signal Controller (\texttt{rl/})}
A DQN agent trained inside SUMO using TraCI. Inputs consist of lane halting
counts, phase state, and time remaining. Output actions adjust green phase
durations dynamically.

% ==========================================
\section{Computer Vision Pipeline}
% ==========================================

\subsection{Stage 1: Lane Detection}
File: \texttt{cv/lane\_detector.py}

\textbf{Purpose:} Compute landmark points \texttt{top}, \texttt{left},
\texttt{right} defining the trapezoidal Region of Interest (ROI).

\subsubsection{Algorithm}
\begin{enumerate}
  \item Rotate each frame 90$^\circ$ CW; resize to $800\times600$.
  \item Greyscale conversion + Canny edge detection (thresholds 50, 150).
  \item Sample Hough Lines every 10th frame
        (\texttt{minLineLength=100}, \texttt{maxLineGap=10}).
  \item Classify line segments by slope $m$:
        \begin{itemize}
          \item $m < -0.4$ $\Rightarrow$ \textbf{blue} (left boundary)
          \item $m > 0.3$  $\Rightarrow$ \textbf{green} (right boundary)
          \item $|m| \le 0.1$ $\Rightarrow$ \textbf{red} (stop line)
        \end{itemize}
  \item Compute temporal averages of lines per colour.
  \item Extend averaged lines across frame boundaries ($y = mx + c$).
  \item Compute intersection via Cramer's rule.
  \item Set \texttt{top}, \texttt{left}, and \texttt{right} landmark points.
\end{enumerate}

\subsubsection{Intersection via Cramer's Rule}
\begin{equation}
\det = A_1 B_2 - A_2 B_1
\end{equation}
\begin{equation}
i_x = \frac{C_1 B_2 - C_2 B_1}{\det},\quad
i_y = \frac{A_1 C_2 - A_2 C_1}{\det}
\end{equation}

\subsection{Stage 2: Background Calibration}
File: \texttt{cv/background\_contour.py}

\textbf{Purpose:} Estimate empty road baseline contour count ($N_{\text{bg}}$)
to subtract static background noise.

\begin{codebox}
\footnotesize
\begin{verbatim}
back_contours = np.min(vehicle_counts) - 20
\end{verbatim}
\end{codebox}

\subsection{Stage 3: Live Vehicle Counting}
File: \texttt{cv/vehicle\_counter.py}

\textbf{Purpose:} Real-time moving vehicle count per approach.

\begin{equation}
\text{count} = \left\lfloor
  \frac{\overline{N}_{\text{contours}} - N_{\text{bg}}}{8}
\right\rfloor
\end{equation}

\subsection{Stage 4: YOLOv8 Counter}
File: \texttt{cv/yolo\_counter.py}

Deep learning alternative evaluating detections for class labels
$\in$ \{car, truck, bus, motorbike\}.
\columnbreak
% ==========================================
\section{SUMO Simulation Environment}
% ==========================================

\subsection{Map Data (\texttt{sumo/maps/})}
Maps include \texttt{SingleTL}, \texttt{intersection}, \texttt{city1},
\texttt{kingcircle}, and \texttt{five\_gardens}.

\subsection{Active SUMO Config}
Configuration file \texttt{sumo/kingcircle.sumocfg} specifies network topology
and vehicle flow parameters.

% ==========================================
\section{Deep Q Network (DQN) Agent}
% ==========================================

\subsection{Gymnasium Environment: \texttt{TrafficLightEnv}}
Implemented in \texttt{rl/train\_dqn.py} and \texttt{rl/test\_dqn.py}.

\subsection{Observation \& Action Spaces}
\begin{equation}
\mathbf{s} = \bigl[h_1, \ldots, h_L,\; \phi_1, \tau_1, \ldots, \phi_K, \tau_K\bigr]
\end{equation}

Action Space: \texttt{Discrete(3)}:
\begin{itemize}
  \item Action 0: Extend green phase (+5s)
  \item Action 1: Switch phase
  \item Action 2: Keep phase
\end{itemize}

\subsection{Reward Function}
\begin{equation}
r_t = -\sum_{\ell \in \mathcal{L}} W_\ell(t)
\end{equation}

% ==========================================
\section{Conclusion}
% ==========================================

Smart Flow demonstrates a modular, scalable architecture combining Computer
Vision and Reinforcement Learning for urban traffic management. With the complete
package refactoring and bug resolution, the project is structured cleanly for
maintainability, further research, and edge deployment.

\end{multicols}
\end{document}
